\newpage

\section{Notation Reference}
\label{app:notation}
\begin{table}[hbt]
\centering
\small
\caption{
Summary of notation. For readability, we omit conditioning variables such as $(x,c)$
when they are clear from context. Throughout the paper, $\pi_\theta$ denotes the student
policy, $\pi_{\mathrm{ref}}$ denotes the reference policy, and $p_T$ denotes the privileged
teacher.
}
\label{tab:notation}
\begin{tabular}{p{0.28\linewidth}p{0.66\linewidth}}
\toprule
\textbf{Notation} & \textbf{Description} \\
\midrule

$x$ & Input prompt. \\
$c$ & Privileged context available only during training. \\
$y=(y_1,\ldots,y_T)$ & Generated reasoning trajectory. \\
$y_t$ & Token at decoding step $t$. \\
$y_{<t}$ & Prefix before token $t$. \\
$h_t=(x,y_{<t})$ & Decoding history or prefix state at step $t$. \\

\midrule

$\pi_\theta(y\mid x)$ & Student policy parameterized by $\theta$. \\
$p_T(y\mid x,c)$ & Privileged teacher policy conditioned on prompt $x$ and context $c$. \\
$\pi_{\mathrm{ref}}(y\mid x)$ & Reference policy, e.g., the initial student or a stop-gradient copy. \\
$\bar{\theta}$ & Stop-gradient copy of the current student parameters. \\

\midrule

$R(y;x,c)$ & Trajectory-level reward in the KL-regularized RL formulation. \\
$\mathcal{J}_{\beta}(\theta)$ & $\beta$-OPSD objective, obtained by using the teacher-to-reference log-ratio reward. \\
$\beta$ & KL regularization coefficient in $\beta$-OPSD. \\
$\beta_k$ & Regularization coefficient used at training step $k$. \\
$\pi_{\beta}^{\star}(y\mid x,c)$ & Closed-form optimal policy of the $\beta$-OPSD objective. \\


\midrule

$z_{\mathrm{ref}}(\cdot\mid h_t)$ & Reference/student logits at prefix $h_t$. \\
$z_T(\cdot\mid h_t,c)$ & Privileged teacher logits at prefix $h_t$. \\
$\tilde p_{\beta_k}(\cdot\mid h_t,c)$ & Local logit-interpolant target at step $k$. \\
$\tilde p_{\beta_k}(y\mid x,c)$ & Autoregressive sequence-level logit-interpolant target. \\
$s_k=\frac{k}{K-1}$ & Normalized training progress. \\
$w_{k}$ & Implementation notation for the teacher weight. \\

\midrule

$\rho_t^{\beta_k}(y)$ & Token-level log-ratio between the student and interpolant target. \\
$G_{t,\gamma}^{\beta_k}(y)$ & Discounted return-to-go from token $t$ using future interpolant log-ratios. \\
$\gamma$ & Discount factor for the return-to-go. \\
$\operatorname{sg}(\cdot)$ & Stop-gradient operator. \\

\midrule

$D_{\mathrm{KL}}(\cdot\|\cdot)$ & Kullback--Leibler divergence. \\
$\mathbb{E}_{y\sim\pi_\theta}$ & Expectation over trajectories sampled from the student policy. \\

\bottomrule
\end{tabular}
\end{table}
\clearpage

\section{Theoretical Proofs}
\label{app:rtg-proofs}
\subsection{Proof of Proposition~\ref{prop:opsd_beta_one}}
\label{app:regularized_rl}
\begin{proof}
Fix a prompt $x$ and privileged context $c$. For readability, we suppress the conditioning
on $(x,c)$ when it is unambiguous. Expanding Eq.~\ref{eq:beta_opsd_objective}, we have
\begin{align}
    \mathcal{J}_{\beta}(\pi_\theta)
    &=
    \mathbb{E}_{y\sim\pi_\theta}
    \left[
    \log
    \frac{
    p_T(y)
    }{
    \pi_{\mathrm{ref}}(y)
    }
    \right]
    -
    \beta
    D_{\mathrm{KL}}
    \left(
    \pi_\theta
    \,\middle\|\,
    \pi_{\mathrm{ref}}
    \right) \\
    &=
    \mathbb{E}_{y\sim\pi_\theta}
    \left[
    \log p_T(y)
    -
    \log \pi_{\mathrm{ref}}(y)
    \right]
    -
    \beta
    \mathbb{E}_{y\sim\pi_\theta}
    \left[
    \log \pi_\theta(y)
    -
    \log \pi_{\mathrm{ref}}(y)
    \right] \\
    &=
    \mathbb{E}_{y\sim\pi_\theta}
    \left[
    \log p_T(y)
    -
    \beta \log \pi_\theta(y)
    +
    (\beta-1)\log \pi_{\mathrm{ref}}(y)
    \right].
    \label{eq:beta_opsd_expanded}
\end{align}
On the other hand,
\begin{align}
    &-
    D_{\mathrm{KL}}
    \left(
    \pi_\theta
    \,\middle\|\,
    p_T
    \right)
    -
    (\beta-1)
    D_{\mathrm{KL}}
    \left(
    \pi_\theta
    \,\middle\|\,
    \pi_{\mathrm{ref}}
    \right) \\
    &=
    -
    \mathbb{E}_{y\sim\pi_\theta}
    \left[
    \log \pi_\theta(y)
    -
    \log p_T(y)
    \right]
    -
    (\beta-1)
    \mathbb{E}_{y\sim\pi_\theta}
    \left[
    \log \pi_\theta(y)
    -
    \log \pi_{\mathrm{ref}}(y)
    \right] \\
    &=
    \mathbb{E}_{y\sim\pi_\theta}
    \left[
    \log p_T(y)
    -
    \beta \log \pi_\theta(y)
    +
    (\beta-1)\log \pi_{\mathrm{ref}}(y)
    \right].
\end{align}
Comparing this expression with Eq.~\ref{eq:beta_opsd_expanded}, we obtain the equivalent
decomposition
\begin{equation}
    \mathcal{J}_{\beta}(\pi_\theta)
    =
    -
    D_{\mathrm{KL}}
    \left(
    \pi_\theta(\cdot\mid x)
    \,\middle\|\,
    p_T(\cdot\mid x,c)
    \right)
    -
    (\beta-1)
    D_{\mathrm{KL}}
    \left(
    \pi_\theta(\cdot\mid x)
    \,\middle\|\,
    \pi_{\mathrm{ref}}(\cdot\mid x)
    \right).
    \label{eq:beta_opsd_decomposition}
\end{equation}
Setting $\beta=1$ cancels the reference-regularization term:
\begin{equation}
    \mathcal{J}_{1}(\pi_\theta)
    =
    -
    D_{\mathrm{KL}}
    \left(
    \pi_\theta(\cdot\mid x)
    \,\middle\|\,
    p_T(\cdot\mid x,c)
    \right).
\end{equation}
Therefore,
\begin{equation}
    \max_{\theta}\mathcal{J}_{1}(\pi_\theta)
    \quad
    \Longleftrightarrow
    \quad
    \min_{\theta}
    D_{\mathrm{KL}}
    \left(
    \pi_\theta(\cdot\mid x)
    \,\middle\|\,
    p_T(\cdot\mid x,c)
    \right),
\end{equation}
which is exactly the standard sequence-level OPSD objective in
Eq.~\ref{eq:standard_opsd_kl}. This proves the proposition.
\end{proof}

\subsection{Proof of Proposition~\ref{prop:geometric_optimal_policy}}
\label{app:geometric-optimal-policy}
\begin{proof}
The $\beta$-OPSD objective:
\begin{equation}
    \mathcal{J}_{\beta}(\pi)
    =
    \mathbb{E}_{y\sim\pi(\cdot\mid x)}
    \left[
    \log
    \frac{
    p_T(y\mid x,c)
    }{
    \pi_{\mathrm{ref}}(y\mid x)
    }
    \right]
    -
    \beta
    D_{\mathrm{KL}}
    \left(
    \pi(\cdot\mid x)
    \,\middle\|\,
    \pi_{\mathrm{ref}}(\cdot\mid x)
    \right).
\end{equation}
Expanding the KL term gives
\begin{align}
    \mathcal{J}_{\beta}(\pi)
    &=
    \mathbb{E}_{y\sim\pi}
    \left[
    \log p_T(y)
    -
    \log \pi_{\mathrm{ref}}(y)
    \right]
    -
    \beta
    \mathbb{E}_{y\sim\pi}
    \left[
    \log \pi(y)
    -
    \log \pi_{\mathrm{ref}}(y)
    \right] \\
    &=
    \mathbb{E}_{y\sim\pi}
    \left[
    \log p_T(y)
    -
    \beta \log \pi(y)
    +
    (\beta-1)\log \pi_{\mathrm{ref}}(y)
    \right].
    \label{eq:beta_opsd_expanded_for_optimal_policy}
\end{align}

Now define
\begin{equation}
    \pi_{\beta}^{\star}(y\mid x,c)
    =
    \frac{
    \pi_{\mathrm{ref}}(y\mid x)^{1-\frac{1}{\beta}}
    p_T(y\mid x,c)^{\frac{1}{\beta}}
    }{
    Z_{\beta}(x,c)
    },
\end{equation}
where
\begin{equation}
    Z_{\beta}(x,c)
    =
    \sum_{y'}
    \pi_{\mathrm{ref}}(y'\mid x)^{1-\frac{1}{\beta}}
    p_T(y'\mid x,c)^{\frac{1}{\beta}}.
\end{equation}
Taking logs, we have
\begin{equation}
    \log \pi_{\beta}^{\star}(y\mid x,c)
    =
    \left(1-\frac{1}{\beta}\right)
    \log \pi_{\mathrm{ref}}(y\mid x)
    +
    \frac{1}{\beta}
    \log p_T(y\mid x,c)
    -
    \log Z_{\beta}(x,c).
\end{equation}
Multiplying both sides by $\beta$ gives
\begin{equation}
    \beta \log \pi_{\beta}^{\star}(y\mid x,c)
    =
    (\beta-1)
    \log \pi_{\mathrm{ref}}(y\mid x)
    +
    \log p_T(y\mid x,c)
    -
    \beta \log Z_{\beta}(x,c).
    \label{eq:beta_log_optimal_policy}
\end{equation}
Substituting Eq.~\ref{eq:beta_log_optimal_policy} into
Eq.~\ref{eq:beta_opsd_expanded_for_optimal_policy}, we obtain
\begin{align}
    \mathcal{J}_{\beta}(\pi)
    &=
    \mathbb{E}_{y\sim\pi}
    \left[
    \beta \log \pi_{\beta}^{\star}(y\mid x,c)
    +
    \beta \log Z_{\beta}(x,c)
    -
    \beta \log \pi(y\mid x)
    \right] \\
    &=
    -\beta
    \mathbb{E}_{y\sim\pi}
    \left[
    \log
    \frac{
    \pi(y\mid x)
    }{
    \pi_{\beta}^{\star}(y\mid x,c)
    }
    \right]
    +
    \beta \log Z_{\beta}(x,c) \\
    &=
    -\beta
    D_{\mathrm{KL}}
    \left(
    \pi(\cdot\mid x)
    \,\middle\|\,
    \pi_{\beta}^{\star}(\cdot\mid x,c)
    \right)
    +
    \beta \log Z_{\beta}(x,c).
    \label{eq:beta_opsd_as_negative_kl}
\end{align}
The second term is independent of $\pi$, while the KL term is nonnegative and is minimized
to zero if and only if
\begin{equation}
    \pi(\cdot\mid x)
    =
    \pi_{\beta}^{\star}(\cdot\mid x,c).
\end{equation}
Therefore, $\pi_{\beta}^{\star}$ is the maximizer of
Eq.~\ref{eq:beta_opsd_objective}. This proves the proposition.
\end{proof}

\subsection{Proof of the Return-to-Go Gradient Estimator}
\label{app:onpolicy-lookahead-gradient}
\begin{proof}
We prove that when $\gamma=1$, the return-to-go estimator in
Eq.~\ref{eq:slide_return_to_go} gives an unbiased estimator of the sequence-level
reverse-KL gradient.

\begin{proposition}[Unbiased sequence-level gradient estimator]
\label{prop:rtg_unbiased_beta_opsd}
Fix a training step $k$ and treat the interpolant target
$\tilde p_{\beta_k}(\cdot\mid x,c)$ as fixed during the current gradient update. Define
\begin{equation}
    \mathcal{L}_{\beta_k}(\theta)
    =
    D_{\mathrm{KL}}
    \left(
    \pi_\theta(\cdot\mid x)
    \,\middle\|\,
    \tilde p_{\beta_k}(\cdot\mid x,c)
    \right),
\end{equation}
and let
\begin{equation}
    \rho_t^{\beta_k}(y)
    =
    \log \pi_\theta(y_t\mid x,y_{<t})
    -
    \log \tilde p_{\beta_k}(y_t\mid x,y_{<t},c).
\end{equation}
For $\gamma=1$, define
\begin{equation}
    G_t^{\beta_k}(y)
    =
    \sum_{s=t}^{T}
    \rho_s^{\beta_k}(y).
\end{equation}
Then
\begin{equation}
    \widehat{g}_{\mathrm{RTG}}(y)
    =
    \sum_{t=1}^{T}
    G_t^{\beta_k}(y)
    \nabla_\theta
    \log \pi_\theta(y_t\mid x,y_{<t})
\end{equation}
satisfies
\begin{equation}
    \mathbb{E}_{y\sim \pi_\theta(\cdot\mid x)}
    \left[
    \widehat{g}_{\mathrm{RTG}}(y)
    \right]
    =
    \nabla_\theta
    \mathcal{L}_{\beta_k}(\theta).
\end{equation}
\end{proposition}

\begin{proof}
For readability, fix $x,c,\beta_k$ and suppress them when clear. By the autoregressive
factorization,
\begin{equation}
    \pi_\theta(y\mid x)
    =
    \prod_{t=1}^{T}
    \pi_\theta(y_t\mid x,y_{<t}),
    \qquad
    \tilde p_{\beta_k}(y\mid x,c)
    =
    \prod_{t=1}^{T}
    \tilde p_{\beta_k}(y_t\mid x,y_{<t},c).
\end{equation}
Therefore,
\begin{equation}
    \mathcal{L}_{\beta_k}(\theta)
    =
    \mathbb{E}_{y\sim\pi_\theta(\cdot\mid x)}
    \left[
    \sum_{t=1}^{T}
    \rho_t^{\beta_k}(y)
    \right].
\end{equation}
Let
\begin{equation}
    R^{\beta_k}(y)
    :=
    \sum_{t=1}^{T}
    \rho_t^{\beta_k}(y),
    \qquad
    g_t(y)
    :=
    \nabla_\theta
    \log \pi_\theta(y_t\mid x,y_{<t}).
\end{equation}
Then
\begin{equation}
    \mathcal{L}_{\beta_k}(\theta)
    =
    \mathbb{E}_{y\sim\pi_\theta}
    \left[
    R^{\beta_k}(y)
    \right].
\end{equation}
Differentiating with the score-function identity gives
\begin{align}
    \nabla_\theta \mathcal{L}_{\beta_k}(\theta)
    &=
    \mathbb{E}_{y\sim\pi_\theta}
    \left[
    R^{\beta_k}(y)
    \nabla_\theta \log \pi_\theta(y\mid x)
    \right]
    +
    \mathbb{E}_{y\sim\pi_\theta}
    \left[
    \nabla_\theta R^{\beta_k}(y)
    \right].
\end{align}
Since
\begin{equation}
    \nabla_\theta \log \pi_\theta(y\mid x)
    =
    \sum_{t=1}^{T}
    g_t(y),
\end{equation}
and the interpolant target $\tilde p_{\beta_k}$ is fixed during the update,
\begin{equation}
    \nabla_\theta R^{\beta_k}(y)
    =
    \sum_{t=1}^{T}
    g_t(y).
\end{equation}
The expected score is zero:
\begin{equation}
    \mathbb{E}_{y\sim\pi_\theta}
    \left[
    \sum_{t=1}^{T}
    g_t(y)
    \right]
    =
    0.
\end{equation}
Thus,
\begin{align}
    \nabla_\theta \mathcal{L}_{\beta_k}(\theta)
    &=
    \mathbb{E}_{y\sim\pi_\theta}
    \left[
    R^{\beta_k}(y)
    \sum_{t=1}^{T}
    g_t(y)
    \right] \\
    &=
    \mathbb{E}_{y\sim\pi_\theta}
    \left[
    \sum_{t=1}^{T}
    g_t(y)
    \sum_{s=1}^{T}
    \rho_s^{\beta_k}(y)
    \right].
    \label{eq:full_return_gradient_beta}
\end{align}

We now remove the past terms by causality. For $s<t$, the quantity
$\rho_s^{\beta_k}(y)$ depends only on earlier tokens and is fixed after conditioning on the
prefix $(x,y_{<t})$. Meanwhile,
\begin{align}
    \mathbb{E}_{y_t\sim\pi_\theta(\cdot\mid x,y_{<t})}
    \left[
    g_t(y)
    \mid x,y_{<t}
    \right]
    &=
    \sum_{v\in\mathcal{V}}
    \pi_\theta(v\mid x,y_{<t})
    \nabla_\theta
    \log \pi_\theta(v\mid x,y_{<t}) \\
    &=
    \sum_{v\in\mathcal{V}}
    \nabla_\theta
    \pi_\theta(v\mid x,y_{<t}) \\
    &=
    \nabla_\theta
    \sum_{v\in\mathcal{V}}
    \pi_\theta(v\mid x,y_{<t}) \\
    &=
    0.
\end{align}
Therefore,
\begin{equation}
    \mathbb{E}_{y\sim\pi_\theta}
    \left[
    g_t(y)
    \sum_{s<t}
    \rho_s^{\beta_k}(y)
    \right]
    =
    0.
\end{equation}
Applying this to Eq.~\ref{eq:full_return_gradient_beta}, only the current and future terms
remain:
\begin{align}
    \nabla_\theta \mathcal{L}_{\beta_k}(\theta)
    &=
    \mathbb{E}_{y\sim\pi_\theta}
    \left[
    \sum_{t=1}^{T}
    g_t(y)
    \sum_{s=t}^{T}
    \rho_s^{\beta_k}(y)
    \right] \\
    &=
    \mathbb{E}_{y\sim\pi_\theta}
    \left[
    \sum_{t=1}^{T}
    G_t^{\beta_k}(y)
    g_t(y)
    \right].
\end{align}
This proves that the return-to-go estimator is unbiased for the sequence-level gradient.
\end{proof}

Finally, the practical surrogate in Eq.~\ref{eq:slide_loss} with $\gamma=1$ has gradient
\begin{equation}
    \nabla_\theta \mathcal{L}_{\ours{}}(\theta;y)
    =
    \frac{1}{T}
    \sum_{t=1}^{T}
    G_t^{\beta_k}(y)
    \nabla_\theta
    \log \pi_\theta(y_t\mid x,y_{<t}),
\end{equation}
where the return weights are detached. Hence, its expectation equals
$\frac{1}{T}\nabla_\theta \mathcal{L}_{\beta_k}(\theta)$. The factor $1/T$ is only a
normalization constant. For $\gamma<1$, the same loss becomes a discounted practical
approximation that reduces the magnitude and variance of long-horizon returns, but the
exact unbiasedness statement corresponds to $\gamma=1$.
\end{proof}


\section{Experimental Details}
\label{sec:exp_details}
\subsection{Training Details}
\label{sec:training_detail}
\paragraph{Shared training configuration.}
Vanilla OPSD and \ours{} use the same model, optimization settings, rollout configuration,
and computational budget. We train Qwen3-1.7B on four NVIDIA RTX A6000 GPUs using
four-way data parallelism with \texttt{accelerate}. Each training rank runs a colocated
single-GPU vLLM instance for on-policy generation. Training is run for a maximum of
$200$ optimization steps with an eight-hour wall-time limit.

\paragraph{Optimization and parameter-efficient tuning.}
We use a learning rate of $5\times10^{-6}$ and clip the gradient norm at $0.1$.
The per-device batch size is $1$, with $8$ gradient-accumulation steps, giving an effective
batch size of
\begin{equation}
    4\ \text{GPUs}
    \times 1\ \text{sample per GPU}
    \times 8\ \text{accumulation steps}
    = 32.
\end{equation}
Gradient checkpointing is enabled. All models are trained in bfloat16 precision using
FlashAttention-2. We use LoRA with rank $r=64$ and scaling parameter
$\alpha_{\mathrm{LoRA}}=128$. LoRA adapters are applied to
\texttt{q\_proj}, \texttt{k\_proj}, \texttt{v\_proj}, \texttt{o\_proj},
\texttt{gate\_proj}, \texttt{up\_proj}, and \texttt{down\_proj}.

\paragraph{On-policy generation.}
Training trajectories are generated with a maximum sequence length of $20{,}000$ tokens
and a maximum completion length of $1{,}024$ tokens. We sample with temperature $1.1$,
top-$p=0.95$, and top-$k=20$. Generation uses vLLM in colocated mode with tensor-parallel
size $1$ and GPU-memory utilization set to $0.5$. Unless otherwise stated, these settings
are shared by Vanilla OPSD and \ours{}, ensuring that their comparison isolates the effect of
the distillation objective and Look-Ahead credit assignment.

\subsection{Evaluation Details}
\label{sec:eval_detail}
\paragraph{Evaluation pipeline.}
We evaluate the Qwen3-1.7B, Qwen3-4B and Qwen3-8B checkpoints using their corresponding
LoRA adapters loaded on top of the original Qwen3 base models.
We evaluate on AIME 2024, AIME 2025, and HMMT 2025.

\paragraph{Decoding configuration.}
We use stochastic decoding with thinking mode enabled. For each problem, we generate
$k=12$ independent candidate solutions using temperature $0.6$, top-$p$ sampling with
$p=0.95$, and top-$k=50$. The evaluation batch size is $32$. We set
\texttt{max\_new\_tokens=0}, corresponding to no separate completion-length cap in the
evaluation script, and use the default maximum model length of $40{,}960$ tokens in
thinking mode. Unless explicitly changed for an ablation, all reported results use this
same decoding configuration.

\paragraph{Scoring and metrics.}
We report
\emph{pass@12}, the fraction of problems for which at least one of the 12 sampled
solutions is correct, and \emph{avg@12}, the average correctness across all 12 sampled
solutions. The evaluation code returns both metrics as raw fractions; throughout the paper,
we multiply them by $100$ and report percentage points.

\paragraph{Inference configuration.}
Evaluation uses tensor-parallel size $1$, GPU memory utilization $0.9$, and a single
evaluation repeat. LoRA checkpoints are evaluated without modifying the underlying base
model weights. Checkpoint sweeps evaluate training steps
$\{50,75,100,200\}$, with one Slurm evaluation job submitted for each
checkpoint--dataset pair.

\section{Sampling from Logit Interpolants}
\label{app:proposal_mixture}
In addition to student on-policy sampling, we evaluate a guided sampling variant that
constructs trajectories from a mixture of the student and privileged teacher. For proposal-mixture weight $\eta\in[0,1]$, we define the token-level proposal distribution
\begin{equation}
    m_{\bar{\theta},\eta}(\cdot\mid h_t,c)
    =
    (1-\eta)\,
    \pi_{\bar{\theta}}(\cdot\mid h_t)
    +
    \eta\,
    p_T(\cdot\mid h_t,c),
    \label{eq:student_teacher_mixed_sampler}
\end{equation}
where $\pi_{\bar{\theta}}$ is a stop-gradient copy of the current student. Thus,
$\eta=0$ recovers student on-policy sampling, whereas $\eta=1$ samples entirely from
the privileged teacher. The resulting autoregressive proposal is
\begin{equation}
    m_{\bar{\theta},\eta}(y\mid x,c)
    =
    \prod_{t=1}^{T}
    m_{\bar{\theta},\eta}(y_t\mid x,y_{<t},c).
    \label{eq:student_teacher_mixed_sequence}
\end{equation}

Because the objective in Eq.~\ref{eq:standard_opsd_kl} is defined under the student trajectory
distribution $\pi_\theta$, sampling from $m_{\bar{\theta},\eta}$ introduces a distribution
mismatch. We correct this mismatch using per-decision importance sampling. For a trajectory
$y\sim m_{\bar{\theta},\eta}(\cdot\mid x,c)$, we define
\begin{equation}
    w_t(y)
    =
    \frac{
    \pi_\theta(y_t\mid x,y_{<t})
    }{
    m_{\bar{\theta},\eta}(y_t\mid x,y_{<t},c)
    },
    \qquad
    W_t(y)
    =
    \prod_{i=1}^{t} w_i(y).
    \label{eq:mixed_sampling_importance_ratio}
\end{equation}
The importance-weighted Look-Ahead return is
\begin{equation}
    G_{t,\gamma}^{\eta,\mathrm{mix}}(y)
    =
    \sum_{s=t}^{T}
    \gamma^{s-t}
    W_s(y)\rho_s^\tau(y),
    \label{eq:mixed_sampling_lookahead_return}
\end{equation}
where $\rho_s^\tau$ measures the student mismatch with the logit-interpolant target. We optimize the corresponding
surrogate
\begin{equation}
    \mathcal{L}_{\mathrm{mix\text{-}\ours}}(\theta;y)
    =
    \frac{1}{T}
    \sum_{t=1}^{T}
    \operatorname{sg}
    \left(
    G_{t,\gamma}^{\tau,\mathrm{mix}}(y)
    \right)
    \log \pi_\theta(y_t\mid x,y_{<t}).
    \label{eq:mixed_sampling_lad_loss}
\end{equation}
We note that this variant is inherently off-policy and requires combining the student and
teacher distributions at every decoding step. At the time of writing, such dual-model sampling is \textbf{not} directly supported by standard vLLM decoding pipelines, resulting in additional (nearly double)
implementation complexity and generation overhead. Although mixed student--teacher
sampling yields comparable improvements over vanilla OPSD, its reduced
sampling efficiency makes it less attractive in practice. 

\paragraph{Results}
We evaluate fixed proposal-mixture weights
$\eta\in\{0.2,0.5,0.8\}$ and linear schedules that either increase or decrease $\eta$
during training. Here, $\eta$ denotes the teacher weight in
Eq.~\ref{eq:student_teacher_mixed_sampler}, while $1-\eta$ denotes the student weight.
All mixed-sampling variants use \ours{} with discount factor $\gamma=0.99$ and
per-decision importance sampling (IS).

\begin{table*}[htb]
\centering
\scriptsize
\caption{
Mixed student--teacher sampling results. All non-baseline variants use
Sampling from Logit Interpolants with discount factor $\gamma=0.99$ and
importance sampling (IS). Here, $\eta$ denotes the teacher weight in the
proposal distribution. We report pass@12 and avg@12 in percentage points.
}
\label{tab:mixed_sampling_results}
\setlength{\tabcolsep}{4pt}
\renewcommand{\arraystretch}{1.06}

\begin{tabular}{lccccccc}
\toprule
Proposal schedule
& Step
& \multicolumn{2}{c}{AIME 2024}
& \multicolumn{2}{c}{AIME 2025}
& \multicolumn{2}{c}{HMMT 2025} \\
\cmidrule(lr){3-4}
\cmidrule(lr){5-6}
\cmidrule(lr){7-8}
&
& pass@12
& avg@12
& pass@12
& avg@12
& pass@12
& avg@12 \\
\midrule

Vanilla OPSD
& 200
& 76.67
& 46.11
& 63.33
& 35.56
& 43.33
& 15.83 \\

\midrule
\multicolumn{8}{l}{\textit{Fixed proposal: $\eta=0.2$}} \\
\quad Mixed Sampling
& 50
& 76.67
& 50.83
& 60.00
& 37.50
& 40.00
& 15.83 \\
\quad Mixed Sampling
& 75
& 80.00
& 50.00
& 66.67
& 38.06
& 40.00
& 14.44 \\
\quad Mixed Sampling
& 100
& 76.67
& 52.78
& 53.33
& 35.83
& 33.33
& 13.61 \\
\quad Mixed Sampling
& 200
& 73.33
& 51.94
& 63.33
& 37.22
& 43.33
& 16.11 \\

\midrule
\multicolumn{8}{l}{\textit{Fixed proposal: $\eta=0.5$}} \\
\quad Mixed Sampling
& 50
& 80.00
& 50.83
& 66.67
& 36.39
& 43.33
& 16.11 \\
\quad Mixed Sampling
& 75
& 80.00
& 52.22
& 66.67
& 37.78
& 46.67
& 16.11 \\
\quad Mixed Sampling
& 100
& 76.67
& \textbf{54.17}
& 63.33
& 39.44
& 46.67
& 17.22 \\
\quad Mixed Sampling
& 200
& 76.67
& 53.33
& 63.33
& 38.06
& 36.67
& 15.83 \\

\midrule
\multicolumn{8}{l}{\textit{Fixed proposal: $\eta=0.8$}} \\
\quad Mixed Sampling
& 50
& 80.00
& 50.83
& 60.00
& 36.94
& 43.33
& 15.83 \\
\quad Mixed Sampling
& 75
& 80.00
& 51.39
& 70.00
& \textbf{41.11}
& 46.67
& 16.11 \\
\quad Mixed Sampling
& 100
& 73.33
& 51.94
& 60.00
& 36.11
& 40.00
& 15.83 \\
\quad Mixed Sampling
& 200
& 73.33
& 51.39
& 66.67
& 38.61
& 36.67
& 13.89 \\

\midrule
\multicolumn{8}{l}{
\textit{Linear proposal schedule: $\eta:0.2\rightarrow0.8$}
} \\
\quad Mixed Sampling
& 50
& 80.00
& 50.83
& 70.00
& 38.06
& 36.67
& 13.61 \\
\quad Mixed Sampling
& 75
& 76.67
& 50.56
& 70.00
& 36.39
& \textbf{53.33}
& 18.06 \\
\quad Mixed Sampling
& 100
& 73.33
& 51.94
& 66.67
& 36.67
& 36.67
& 16.39 \\
\quad Mixed Sampling
& 200
& 76.67
& 50.56
& 60.00
& \textbf{41.11}
& 43.33
& 16.94 \\

\midrule
\multicolumn{8}{l}{
\textit{Linear proposal schedule: $\eta:0.8\rightarrow0.2$}
} \\
\quad Mixed Sampling
& 50
& 76.67
& 50.56
& \textbf{76.67}
& 36.39
& 30.00
& 13.61 \\
\quad Mixed Sampling
& 75
& 73.33
& 47.78
& 73.33
& 36.39
& 50.00
& \textbf{18.89} \\
\quad Mixed Sampling
& 100
& 80.00
& 50.56
& 60.00
& 38.61
& 36.67
& 15.83 \\
\quad Mixed Sampling
& 200
& 76.67
& 52.78
& 60.00
& 39.72
& 36.67
& 15.56 \\

\bottomrule
\end{tabular}
\end{table*}

Table~\ref{tab:mixed_sampling_results} shows that mixed student--teacher sampling can
produce improvements comparable to those of the main on-policy \ours{} method. At 200
steps, every evaluated schedule improves avg@12 over Vanilla OPSD on both AIME
benchmarks. Across all checkpoints, the best observed gains over Vanilla OPSD are $8.06$
points on AIME 2024, $5.55$ points on AIME 2025, and $3.06$ points on HMMT 2025.
However, no single proposal schedule dominates across all benchmarks, indicating that the
off-policy variant is sensitive to both the mixture coefficient and the training checkpoint.
Given its additional dual-model decoding cost and the variance introduced by importance
sampling, we use student on-policy sampling as the default implementation of \ours{}.